\chapter{总结与收获}

本次课程设计在华为云上完整走通了「容器化 $\to$ K8s 编排 $\to$ 大数据计算 $\to$ 运维监控
$\to$ 持续交付 $\to$ 分布式训练」的云原生全链路，从把概念落到了真实生产环境。

\textbf{对云计算技术的新认识。} 其一，\emph{声明式与不可变带来的确定性}。整个平台用 YAML
声明期望状态，\lstinline!kubectl apply! 后由控制器持续调谐，副本数、镜像版本、配置都可
版本化、可复现；多阶段构建产出的不可变镜像（后端仅 202\,MB）配合 SWR 私有仓库，使
「一次构建、处处运行」成为现实。其二，\emph{关注点分离}。ConfigMap/Secret 把配置与镜像
解耦，PVC 把数据生命周期与 Pod 解耦（删除并重建 Redis Pod 后 \lstinline!testkey! 仍在），
Service/ELB 把访问入口与具体 Pod 解耦，每一层都可独立演进。其三，\emph{弹性即成本}。HPA
让后端副本在 1$\sim$4 之间随真实 CPU 负载自动伸缩，低峰省钱、高峰保障 SLA，直接体现
「按需付费」。其四，\emph{分布式的代价与收益}。Spark 与 PyTorch DDP 的实验让我直观看到
并行加速并非线性——通信、序列化与数据规模决定了 Amdahl 定律中的串行瓶颈。

\textbf{主要挑战与解决。} 实践中遇到并解决了一系列真实工程问题：① 国内网络无法直连
Docker Hub / ghcr.io / quay.io，通过 daocloud 镜像代理拉取后重打 Tag 推送到 SWR 解决；
② Docker~29.x 默认 OCI manifest 与 SWR 不兼容，构建加 \lstinline!--provenance=false!；
③ 集群默认未启用 metrics-server 导致 HPA 无数据，手动部署并加 \lstinline!--kubelet-insecure-tls!；
④ 2 节点（各 2\,vCPU/8\,GiB）资源紧张，运行 Spark/训练作业前先把 Web 应用副本缩容、
并用 \lstinline!coreRequest! 把 Spark Executor 的 K8s CPU 请求与逻辑核心解耦，规避
\lstinline!FailedScheduling!；⑤ ConfigMap 的 \lstinline!subPath! 挂载不会热更新，
通过删除 Pod 重建加载新配置。这些"坑"恰恰是云上工程最有价值的经验。

总体而言，本设计不仅完成了全部规定任务与三个附加专题，更让我体会到云原生体系中
\emph{抽象层层叠加、控制器持续调谐、资源弹性供给}的设计哲学，对今后从事云计算与
大数据相关开发打下了扎实基础。
